\section{Lecture 13 (October 30th)}
\begin{thm}
Let $h'(z_0)=\mathrm{sup}\{|\phi '(z_0)|\;|\; \phi \in \mathrm{\Sigma} \}$ then $h(\mathrm{\Omega} )=U$ and $h\in \mathrm{\Sigma} $. In addition, $h(z_0)=0$. In addition, if we add the condition that $h'(z_0)>0$, such $h$ is uniquely determined. 
\end{thm}
\vspace{2ex}
\begin{proof}
If $h'(z_0)=re^{i\theta }$, then $f(z)=e^{-i\theta }h(z)$ satisfies $f\in \mathrm{\Sigma} $, $f(\mathrm{\Omega} )=U$, $f(z_0)=0$ and $f'(z_0)=r>0$. Let's prove uniqueness. If there exists $f$ and $g$ satisfying all conditions, then define
\[h(z)=f\Big(g^{-1}(z)\Big)\]
and $h:U\rightarrow U$ is 1-1 analytic and onto.
\[h'(0)=\dfrac{f'(z_0)}{g'(z_0)}>0\]
and $h(0)=0$. This implies that $h(z)=0$ and that $f=g$. 
\end{proof}
\vspace{2ex}
\begin{thm}
Let $D$ be simply connected and $D\ne {\bm C}$. Then there is $\phi :D\leftarrow U$ that is 1-1 analytic and onto. Recall that if $\phi $ is 1-1, $\phi '\ne 0$ and is conformal. This tells us that there exists $\psi =\phi ^{-1}:U\rightarrow D$ that is still 1-1 analytic and onto, thus conformal. Define $A(D)$ be the set of all analytic functions on $D$. Define $\mathrm{\Phi} :A(D)\rightarrow A(U)$ by $\mathrm{\Phi} (f)=f\circ \psi $. Then, $\mathrm{\Phi} $ becomes a ring isomorphism. 
\end{thm}
\vspace{2ex}
\begin{rmk}
Consider ${\bm C}^2={\bm C}\times {\bm C}$. The simplest domains in ${\bm C}^2$ are $B_{2}$ and $U^2$.
\[B_{2}=\{z=(z_1,z_2)\in {\bm C}^2 \;|\;||z||<1 \}\qquad\&\qquad U^2=U\times U\]
each called a bidisc and a unit ball. The norm $||z||=\langle z,z\rangle ^{1/2}$ is defined as
\[\langle z,z\rangle ^{1/2}=\langle (z_1,z_2),(z_1,z_2)\rangle ^{1/2}=\Big(z_1\bar{z}_{1}+z_{2}\bar{z}_{2}\Big)^2=\Big(|z_1|^2+|z_2|^2\Big)^{1/2}\]
\end{rmk}
\vspace{2ex}
\begin{thm}
(Residue theorem) If $f$ is analytic inside and on a POSCC $C$ except for isolated singularities at $z_1,\ldots,z_{n}$ inside $C$, then
\[\int _{C}f(z)\,dz=2\pi i\sum ^{n}_{k=1}\mathop{\mathrm{Res}}_{z=z_{k}}f(z)\]
If $f(z)=\sum ^{\infty }_{n=-\infty }a_{n}(z-z_0)^{n}$ in $0<|z-z_{k}|<\delta $ (Laurent series of $f$ at $z_k$), then 
\[a_{-1}=\mathop{\mathrm{Res}}_{z=z_{k}}f(z)\]
If $\lim _{z\rightarrow z_0}(z-z_0)^{k}f(z_0)=A\ne 0$, we say that $f$ has a pole of order $k$ at $z_0$. If $f$ has a simple pole (order 1) at $z_0$, then $\mathop{\mathrm{Res}}_{z=z_0}f(z)=\lim _{z\rightarrow z_0}(z-z_0)f(z)$. If 
\[f(z)=\dfrac{p(z)}{q(z)}\]
($p,q$ analytic at $z_0$) has a simple pole at $z_0$ ($p(z_0)\ne 0 $, $q(z_0)=0$, and $q'(z_0)\ne 0$), then
\[\mathop{\mathrm{Res}}_{z=z_0}f(z)=\lim _{z\rightarrow z_0}(z-z_0)\dfrac{p(z)}{q(z)}=\lim _{z\rightarrow z_0}\dfrac{p(z)}{\dfrac{q(z)-q(z_0)}{z-z_0}}=\dfrac{p(z_0)}{q'(z_0)}\]
\end{thm}
\vspace{2ex}
\begin{ex}
Evaluate for $0<\alpha <1$,
\[\int ^{\infty }_{-\infty }\dfrac{e^{\alpha x}}{1+e^{x}}\,dx\]
\end{ex}
\vspace{2ex}
\begin{proof}
We need to take a rectangular contour $C_{R}$ from $-R$ to $R$ and height $2\pi i$ notice that the above function has a residue at $\pi i$. Take 
\[f(z)=\dfrac{\alpha z}{1+e^{z}}\]
For every $R>0$,
\[\int _{C_{R}}f(z)\,dz=2\pi i\mathop{\mathrm{Res}}_{z=\pi i}f(z)=2\pi i\dfrac{e^{\alpha \pi i}}{e^{\pi i}}=-2\pi ie^{\alpha \pi i}\]
On the other hand,
\[\int _{C_{R}}f(z)\,dz=\int ^{R}_{-R}\dfrac{e^{\alpha x}}{1+e^{x}}\,dx+\int ^{2\pi }_{0}\dfrac{e^{\alpha (R+iy)}}{1+e^{R+iy}}i\,dy+\int ^{-R}_{R}\dfrac{e^{\alpha (x+2\pi i)}}{1+e^{x+2\pi i}}\,dx+\int ^{0}_{2\pi }\dfrac{e^{\alpha (-R+iy)}}{1+e^{-R+iy}}i\,dy\]
The second term goes to zero as $R\rightarrow \infty $ and 
\[-2\pi ie^{\alpha \pi i}=\int ^{\infty }_{-\infty }\dfrac{e^{\alpha x}}{1+e^{x}}\,dx\Big[1-e^{2\alpha \pi i}\Big]\]
resulting in
\[\int ^{\infty }_{-\infty }\dfrac{e^{\alpha x}}{1+e^{x}}\,dx=-\dfrac{2\pi i}{1-e^{2\alpha \pi i}}=\dfrac{\pi }{\dfrac{e^{\alpha \pi i}-e^{-\alpha \pi i}}{2i}}=\dfrac{\pi }{\sin \alpha \pi }\]
If we further put $t=e^{x}$, $dx=e^{-x}$,  and $dt/t=dx$. We therefore have
\[\int ^{\infty }_{0}\dfrac{t^{\alpha -1}}{1+t}\,dt=\dfrac{\pi }{\sin \alpha \pi }\]

\end{proof}
\vspace{2ex}

